\documentclass[12pt, a4paper]{article}

% --- Paquetages ---
\usepackage[utf8]{inputenc}
\usepackage[T1]{fontenc}
\usepackage[french]{babel}
\usepackage{graphicx}
\usepackage{geometry}
\geometry{hmargin=2.5cm,vmargin=2.5cm}
\usepackage{hyperref}
\usepackage{titlesec}
\usepackage{listings}
\usepackage{xcolor}
\usepackage{caption}
\usepackage{subcaption}
\usepackage{enumitem}
\usepackage{booktabs}
\usepackage{float}
\usepackage{url}
\usepackage{microtype}

% --- Chemins/URLs en police à chasse fixe, avec retour à la ligne sur "/" et "_" ---
\urlstyle{tt}
\newcommand{\code}[1]{\path{#1}}

% --- Évite les débordements de justification (Overfull \hbox) en autorisant
%     un léger étirement supplémentaire des espaces en dernier recours ---
\setlength{\emergencystretch}{3em}

% --- Paramétrage du code source ---
\lstset{
    basicstyle=\ttfamily\small,
    breaklines=true,
    backgroundcolor=\color{gray!10},
    frame=single,
    keywordstyle=\color{blue},
    commentstyle=\color{green!50!black},
    stringstyle=\color{red}
}

% --- Macro pour emplacement de figure ---
\newcommand{\figureplaceholder}[2]{%
    \fbox{\parbox{0.92\textwidth}{\centering\vspace{3.5cm}%
    \textit{#1}\\[0.5cm]\small #2\vspace{3.5cm}}}%
}

\begin{document}

% --- Page de Garde ---
\begin{titlepage}
	\newcommand{\HRule}{\rule{\linewidth}{0.5mm}}

	\center

	\textsc{\LARGE Université du Québec à Chicoutimi}\\[1.5cm]
	\textsc{\Large Département d'informatique et de mathématique (DIM)}\\[0.5cm]

	\HRule\\[0.4cm]
	{\huge\bfseries 8INF950 --- Sujets spéciaux\\[0.2cm]
	Vibe Coding \& Qualité logicielle en santé\\[0.3cm]
	\large Compte rendu final}\\[0.4cm]
	\HRule\\[1.5cm]

	\begin{center}
		\large
		\textit{Étudiant}\\
		TOÏGO Nathan
	\end{center}

	\vspace{1.5cm}
	{\large\textit{Remis à}}\\
	Hamid \textsc{Mcheick}\\
	Professeur

	\vfill
	{\large\today}
	\vfill
\end{titlepage}

\newpage
\tableofcontents
\vspace{1cm}
\listoffigures
\newpage

% --- Corps du compte rendu ---

\section*{Résumé}
\addcontentsline{toc}{section}{Résumé}

Ce projet compare l'impact du \emph{vibe coding} sur la qualité logicielle en santé.
Quatre modèles Azure AI Foundry (\texttt{gpt-5.4-nano}, \texttt{gpt-5.3-codex},
\texttt{gpt-5.6-sol}, \texttt{gpt-5.6-terra}) ont généré, à partir du même prompt,
une API FastAPI de gestion de rendez-vous médicaux, sans révision humaine. Chaque sortie
a été enregistrée dans un \code{GENERATED_DRAFT.md}, puis matérialisée en fichiers
exécutables via \code{extract_code_blocks.py}. Les quatre implémentations ont été
analysées par SonarQube (\code{http://localhost:9000}). Les résultats montrent des
écarts marqués~: \texttt{gpt-5.6-terra} obtient le meilleur profil (4~violations,
dette 20~min), tandis que \texttt{gpt-5.4-nano} est le seul à présenter
5~vulnérabilités (note sécurité~C). Une \textbf{analyse étendue} complémentaire (lint,
couverture de tests $\geq 80\,\%$, audit de dépendances) a ensuite été appliquée aux
\emph{mêmes} quatre projets~: aucun ne satisfait ce quality gate plus strict, principalement
en raison d'une couverture de tests quasi nulle et de dépendances vulnérables non détectées
par l'analyse SonarQube seule.

\section{Introduction}

\subsection{Contexte}

Les modèles générateurs de code transforment les pratiques de développement. Le
\emph{vibe coding} consiste à formuler des intentions en langage naturel et à accepter
largement les suggestions produites, sans revue systématique. En santé, cela pose des
enjeux de confidentialité (Loi~25), de fiabilité des règles métier et de traçabilité.

Le cours 8INF950 (Été 2026, UQAC) évalue ces pratiques via la norme \textbf{ISO/IEC~25010}
et des outils de mesure objectifs, principalement \textbf{SonarQube}.

\subsection{Problématique}

\emph{Quel modèle générateur de code produit la meilleure qualité logicielle en vibe
coding pur, à spécification identique, et cette qualité résiste-t-elle à une
\textbf{analyse étendue} (lint, couverture de tests, audit de dépendances) allant au-delà
du scan \textbf{SonarQube}~?}

\subsection{Portée}
Ce projet contient : une spécification commune des attentes fonctionnelles
(\code{docs/cahier-des-charges.md}), quatre projets générés par LLM externe (Azure AI
Foundry) à partir de cette spécification, ainsi que deux analyses complémentaires appliquées
à ces mêmes projets~: une \textbf{analyse SonarQube} (statique) et une \textbf{analyse
étendue} (lint, tests, audit de dépendances). Il n'est pas inclus d'interface graphique (UI)
pour naviguer et utiliser les solutions produites, la qualité du code ne pouvant être
analysée uniquement sur la base du code et non pas de la pertinence ou des choix esthétiques
des interfaces.

\subsection{Choix des modèles génératifs}

Les quatre modèles comparés (\texttt{gpt-5.4-nano}, \texttt{gpt-5.3-codex},
\texttt{gpt-5.6-sol}, \texttt{gpt-5.6-terra}) appartiennent tous à la famille
\textbf{OpenAI}, déployés via \textbf{Azure AI Foundry}. Ce choix n'est pas un parti pris
méthodologique visant à écarter d'autres fournisseurs, mais une \textbf{contrainte
pratique}~: l'accès aux ressources de calcul repose sur l'abonnement
\emph{Azure for Students}, qui n'ouvre l'accès qu'aux déploiements OpenAI disponibles sur
Azure AI Foundry (à l'exclusion, par exemple, de modèles Anthropic ou Google, non proposés
dans ce cadre). La comparaison porte donc sur quatre \emph{générations/tailles} différentes
d'un même fournisseur plutôt que sur plusieurs fournisseurs concurrents, ce qui reste
cohérent avec la problématique du projet, centrée sur la variabilité de la qualité du code
généré et non sur un comparatif inter-fournisseurs.

\begin{figure}[H]
    \centering
    \begin{figure}[H]
    \centering
    \begin{lstlisting}[basicstyle=\ttfamily\scriptsize, frame=single, backgroundcolor=\color{gray!8}]
8INF950ETE2026/
|-- docs/ ..................... cahier des charges, grille ISO 25010, etat de l'art
|-- scripts/ .................. generate_with_llm, extract_code_blocks, compare_metrics
`-- case-study/
    |-- docker-compose.sonarqube.yml
    |-- variant-a-nano/ ......... gpt-5.4-nano 
    |-- variant-a-codex/ ........ gpt-5.3-codex
    |-- variant-a-sol/ .......... gpt-5.6-sol
    |-- variant-a-terra/ ........ gpt-5.6-terra
    `-- variant-b-hybrid/ ....... config. CI de reference (analyse etendue)
    \end{lstlisting}
    \caption{Organisation du dépôt \texttt{8INF950ETE2026}~: les quatre projets générés
    (\texttt{variant-a-*}) et la configuration de référence de l'analyse étendue
    (\texttt{variant-b-hybrid}), qui n'est pas une cinquième implémentation comparée.}
    \label{fig:repo-overview}
\end{figure}
\end{figure}

\section{Étude du marché}
\label{sec:etude-marche}

\subsection{Vibe coding et qualité logicielle}

Le vibe coding accélère la production initiale, mais la littérature signale des risques~:
hallucinations, absence de tests, incohérences architecturales et vulnérabilités non
détectées. Plutôt que des critères qualitatifs sur les IDE (Cursor, Copilot), cette étude
retient des \textbf{métriques SonarQube} mesurables et reproductibles
(\code{docs/etat-de-lart-outils-ia.md}).

\subsection{Modèles générateurs comparés}

Quatre déploiements Azure AI Foundry ont été testés avec le même prompt
(\code{docs/prompts/redaction-projet-medical.md}) et la commande~:
\begin{lstlisting}[language=bash,caption={Génération reproductible hors IDE.}]
python scripts/generate_with_llm.py --provider azure \
  --project-name "Gestion de rendez-vous medicaux" \
  --output case-study/variant-a-terra/GENERATED_DRAFT.md \
  --code-dir case-study/variant-a-terra
\end{lstlisting}

Le tableau~\ref{tab:sonar-models} synthétise les résultats SonarQube collectés le
1er~août~2026 sur \code{http://localhost:9000}.

\begin{table}[H]
    \centering
    \caption{Comparaison SonarQube des quatre modèles générateurs (analyse SonarQube)}
    \label{tab:sonar-models}
    \small
    \begin{tabular}{@{}lcccc@{}}
        \toprule
        \textbf{Métrique} & \textbf{nano} & \textbf{codex} & \textbf{sol} & \textbf{terra} \\
        \midrule
        Bugs & 0 & 0 & 0 & 0 \\
        Vulnérabilités & 5 & 0 & 0 & 0 \\
        Code smells & 30 & 48 & 30 & 4 \\
        Couverture (\%) & 0,0 & 0,0 & 0,0 & 0,0 \\
        Duplication (\%) & 0,0 & 0,0 & 0,0 & 0,0 \\
        Complexité cyclomatique & 57 & 82 & 53 & 47 \\
        Dette technique (min) & 151 & 242 & 151 & 20 \\
        Security hotspots & 0 & 0 & 0 & 0 \\
        Note fiabilité & A & A & A & A \\
        Note sécurité & C & A & A & A \\
        Note maintenabilité & A & A & A & A \\
        Lignes de code (NCLOC) & 587 & 585 & 650 & 529 \\
        Violations (total) & 35 & 48 & 30 & 4 \\
        \bottomrule
    \end{tabular}
\end{table}

\subsection{Synthèse préliminaire}

\begin{itemize}[leftmargin=*]
    \item \textbf{\texttt{gpt-5.6-terra}}~: meilleur profil global (4~violations, dette 20~min).
    \item \textbf{\texttt{gpt-5.3-codex}}~: complexité la plus élevée (82) et dette maximale (242~min).
    \item \textbf{\texttt{gpt-5.6-sol}}~: base de code la plus volumineuse (650~lignes), profil
    intermédiaire (30~code smells, dette 151~min --- comme nano) mais sans vulnérabilité
    détectée (note sécurité~A).
    \item \textbf{\texttt{gpt-5.4-nano}}~: seul modèle avec 5~vulnérabilités et note sécurité~C.
    \item Couverture à 0\,\% pour tous (absence de \texttt{coverage.xml} --- attendu en vibe coding pur).
\end{itemize}

\begin{figure}[H]
    \centering
    \includegraphics[width=0.95\textwidth]{fig2.png}
    \caption{Tableaux de bord SonarQube des quatre modèles générateurs.}
    \label{fig:sonar-dashboards}
\end{figure}

\section{Approches informatiques}

\subsection{Spécification fonctionnelle commune}

Le cahier des charges (\code{docs/cahier-des-charges.md}) définit les entités Patient,
Médecin et Rendez-vous, avec règles de chevauchement, validation des entrées et endpoint
\code{/health}. Chaque modèle reçoit la même consigne via le prompt standardisé.

\subsection{Chaîne de génération (vibe coding pur)}

La méthodologie reproductible repose sur trois étapes~:

\begin{enumerate}[leftmargin=*]
    \item \textbf{Prompt}~: \code{docs/prompts/redaction-projet-medical.md} injecté dans
    \code{generate_with_llm.py} (Azure OpenAI compatible \code{/openai/v1} ou LLM local).
    \item \textbf{Brouillon}~: le LLM produit un \code{GENERATED_DRAFT.md} balisé avec des
    marqueurs \code{===FILE: chemin/relatif===} suivis de blocs de code.
    \item \textbf{Extraction}~: \code{extract_code_blocks.py} matérialise les fichiers
    sur disque sans révision du contenu généré.
\end{enumerate}

\begin{figure}[H]
    \centering
    \includegraphics[width=\textwidth]{fig3.png}
    \caption{Pipeline de génération et d'analyse des quatre projets.}
    \label{fig:gen-pipeline}
\end{figure}

\subsection{Architecture des projets générés}

Chaque modèle produit une structure légèrement différente (choix typique du vibe coding).
L'exemple \texttt{variant-a-terra} expose une API FastAPI modulaire~:

\begin{lstlisting}[language=Python,caption={Point d'entrée \code{variant-a-terra/app/main.py}.}]
app = FastAPI(title=settings.app_name, version=settings.app_version, lifespan=lifespan)
app.include_router(patients.router)
app.include_router(practitioners.router)
app.include_router(appointments.router)

@app.get("/health", tags=["Supervision"])
def health_check():
    return {"status": "ok", "service": settings.app_name, "version": settings.app_version}
\end{lstlisting}

Les routes couvrent typiquement CRUD patients, praticiens/médecins et rendez-vous, avec
validation Pydantic et persistance SQLite via SQLAlchemy. La documentation Swagger est
disponible sur \code{/docs}.

\begin{figure}[H]
    \centering
    \includegraphics[width=\textwidth]{fig4.png}
    \caption{Swagger de l'application générée par le modèle gpt-5.6-sol sur le port 8001.}
    \label{fig:gen-pipeline}
\end{figure}

\subsection{Analyse étendue (quality gate CI)}
\label{sec:analyse-etendue}

En complément de l'analyse SonarQube (scan statique ponctuel, section~\ref{sec:etude-marche}),
une \textbf{analyse étendue} a été appliquée aux \emph{mêmes quatre projets générés}, en
reproduisant le quality gate attendu avant toute mise en production~: lint, couverture de
tests minimale et audit des dépendances. La configuration de référence de ce pipeline est
formalisée dans \code{case-study/variant-b-hybrid/README.md}~:

\begin{enumerate}[leftmargin=*]
    \item \textbf{Ruff} --- \texttt{ruff check app tests} (style, erreurs courantes).
    \item \textbf{Pytest} --- couverture minimale de 80\,\% (\texttt{--cov-fail-under=80}).
    \item \textbf{pip-audit} --- audit des vulnérabilités connues des dépendances.
\end{enumerate}

Ce pipeline est outillé via \textbf{GitHub Actions}
(\code{.github/workflows/variant-b-hybrid-ci.yml}), déclenché à chaque \texttt{push} ou
\emph{pull request}. Le dossier \code{case-study/variant-b-hybrid} ne contient que la
configuration CI de référence (workflow, dépendances de développement) ainsi qu'un point de
terminaison \code{/health} minimal ayant servi à valider l'outillage avant son application
aux quatre projets~: il ne constitue \textbf{pas} une cinquième implémentation comparée dans
cette étude --- les projets analysés restent ceux de la section précédente.

\subsubsection{Résultats}

Le pipeline a été rejoué localement sur les quatre projets générés via
\texttt{make ci-variants-a} (\code{scripts/run_ci_local.py})~: le
tableau~\ref{tab:analyse-etendue} montre qu'\textbf{aucun} des quatre projets ne satisfait ce
quality gate, principalement en raison d'une couverture de tests insuffisante et de
dépendances vulnérables --- deux angles morts que l'analyse SonarQube seule ne fait pas
apparaître avec la même sévérité (cf. tableau~\ref{tab:sonar-models}, notes de fiabilité et
de maintenabilité majoritairement~A). Ce constat appuie l'hypothèse~H1.

\begin{table}[H]
    \centering
    \caption{Résultats de l'analyse étendue (pipeline CI local, 5~août~2026) sur les quatre
    projets générés. Modèles : nano = \texttt{gpt-5.4-nano}, codex = \texttt{gpt-5.3-codex},
    sol = \texttt{gpt-5.6-sol}, terra = \texttt{gpt-5.6-terra}.}
    \label{tab:analyse-etendue}
    \footnotesize
    \setlength{\tabcolsep}{4pt}
    \begin{tabular}{@{}lcccc@{}}
        \toprule
        \textbf{Métrique} & \textbf{nano} & \textbf{codex} & \textbf{sol} & \textbf{terra} \\
        \midrule
        Ruff (lint) & échec & échec & échec & échec \\
        Couverture Pytest & 10\,\% & 33\,\% & 60\,\% & 66\,\% \\
        Seuil $\geq$ 80\,\% & non & non & non & non \\
        pip-audit & 8 vuln. & 8 vuln. & 8 vuln. & 8 vuln. \\
        Quality gate & FAIL & FAIL & FAIL & FAIL \\
        \bottomrule
    \end{tabular}
\end{table}

\subsection{Outils de mesure}

Les deux analyses (SonarQube et étendue) sont appliquées aux mêmes quatre projets, avec les
mêmes outils et la même configuration pour chaque modèle, afin de garantir l'égalité de
jugement. SonarQube Community est déployé via Docker
(\code{case-study/docker-compose.sonarqube.yml}).

\subsection{Difficultés techniques rencontrées}

\begin{itemize}[leftmargin=*]
    \item \textbf{Extraction de code}~: le regex initial de \code{extract_code_blocks.py}
    ratait les blocs vides (\code{__init__.py}) et \texttt{database.py}~; correction
    via une fence fermante en début de ligne (\texttt{\^{}```}).
    \item \textbf{SonarQube}~: configuration \texttt{sonar.tests} et exclusions selon
    l'emplacement des tests (\code{tests/} vs \code{app/tests/}).
    \item \textbf{Qualité du code généré}~: emails invalides dans les tests
    (\texttt{@example.test}), architectures hétérogènes entre modèles --- illustrant les
    risques documentés dans \code{docs/risques-vibe-coding.md}.
\end{itemize}

\section{Analyse}

\subsection{Hypothèses de recherche}

\begin{description}[leftmargin=3.5cm, style=nextline]
    \item[H1] Une analyse SonarQube seule sous-estime les lacunes qualité (couverture de
    tests, dépendances vulnérables) que révèle une analyse étendue (lint~+~tests~+~audit).
    \item[H2] Certains modèles génèrent davantage de vulnérabilités (ex. nano vs terra).
    \item[H3] La couverture de tests est quasi nulle sans contrainte explicite imposée par
    un quality gate.
    \item[H4] Le temps de génération initial est court, mais le coût de mise en conformité
    avec l'analyse étendue varie selon le modèle.
\end{description}

\subsection{Interprétation des résultats SonarQube}
    

\subsection{Limites}

\begin{itemize}[leftmargin=*]
    \item Une seule étude de cas et un seul prompt par modèle.
    \item Couverture SonarQube non alimentée (pas de \texttt{coverage.xml}).
    \item L'analyse étendue (lint, tests, audit) n'est pas intégrée aux tableaux de bord
    SonarQube~; les deux analyses restent consultées séparément, sans vue unifiée.
\end{itemize}

\section{Résultats}

\subsection{Résultats de l'analyse SonarQube}

Les quatre projets ont été analysés avec succès sur SonarQube Community. Le tableau~\ref{tab:sonar-models} constitue le livrable
principal de comparaison inter-modèles.

\begin{figure}[H]
    \centering
    \includegraphics[width=0.95\textwidth]{fig6.png}
    \caption{Diagramme en barres~: violations, code smells et dette technique par modèle
    (généré par \code{scripts/generate_fig6_chart.py}).}
    \label{fig:comparison-chart}
\end{figure}

\subsection{Résultats de l'analyse étendue}

Le tableau~\ref{tab:analyse-etendue} (section~\ref{sec:analyse-etendue}) confirme qu'aucun
des quatre projets ne satisfait le quality gate étendu, essentiellement à cause de la
couverture de tests et des dépendances vulnérables --- deux angles morts que l'analyse
SonarQube seule ne révèle pas avec la même sévérité.

\section{Conclusion}

\subsection{Bilan}

L'étude a mis en place une chaîne reproductible de vibe coding (prompt $\rightarrow$ brouillon
$\rightarrow$ extraction) et comparé quatre modèles Azure sur une spécification commune, au
moyen de deux analyses complémentaires~: une \textbf{analyse SonarQube} (statique) et une
\textbf{analyse étendue} (lint, couverture de tests, audit de dépendances). Les écarts de
qualité sont significatifs~: \texttt{gpt-5.6-terra} ressort comme le meilleur candidat en
vibe coding pur, \texttt{gpt-5.4-nano} comme le plus risqué côté sécurité --- mais aucun des
quatre modèles ne satisfait le quality gate de l'analyse étendue.

\subsection{Réponse à la problématique}

Le vibe coding pur produit des applications fonctionnelles rapidement, mais la qualité
dépend fortement du modèle choisi. Les métriques SonarQube objectivent ces différences sans
revue humaine, mais elles ne suffisent pas~: l'analyse étendue révèle, sur les \emph{mêmes}
projets, des lacunes qu'un scan statique seul ne détecte pas (couverture de tests quasi
nulle, dépendances vulnérables), faisant échouer les quatre modèles au quality gate combiné.
Une analyse purement SonarQube ne garantit donc pas, à elle seule, la qualité logicielle
attendue en contexte de santé.

\section*{Références}
\addcontentsline{toc}{section}{Références}

\begin{enumerate}[label={[\arabic*]}, leftmargin=*]
    \item ISO/IEC 25010:2023 --- Systems and software engineering --- SQuaRE --- Product quality model.
    \item SonarSource, \emph{SonarQube Documentation}, \url{https://docs.sonarsource.com/sonarqube/}.
    \item FastAPI, \emph{Documentation officielle}, \url{https://fastapi.tiangolo.com/}.
    \item Mcheick, H., \emph{Plan de cours 8INF950 --- Été 2026}, UQAC (\texttt{plan8INF950ETE2026.pdf}).
    \item TOÏGO, N., \emph{État de l'art --- Comparaison SonarQube des modèles IA},
    \code{docs/etat-de-lart-outils-ia.md}.
    \item Commission d'accès à l'information du Québec, \emph{Loi~25}.
\end{enumerate}

\end{document}
